%% =====================================================================
%% sec_model.tex -- Theory chapter: kinematics, dynamics, and the
%% parameter structure of the payload-adaptive model.
%%
%% USAGE:  %% =====================================================================
%% sec_model.tex -- Theory chapter: kinematics, dynamics, and the
%% parameter structure of the payload-adaptive model.
%%
%% USAGE:  %% =====================================================================
%% sec_model.tex -- Theory chapter: kinematics, dynamics, and the
%% parameter structure of the payload-adaptive model.
%%
%% USAGE:  %% =====================================================================
%% sec_model.tex -- Theory chapter: kinematics, dynamics, and the
%% parameter structure of the payload-adaptive model.
%%
%% USAGE:  \input{sec_model}   (place after \section{Related Work};
%%         it supersedes the "Dynamics and parameterization" / "MHE" /
%%         "Inner-loop gain scaling" paragraphs of the current
%%         \section{System Overview and Problem Setup}, and subsumes
%%         sec3_parameters.tex in full.)
%%
%% REQUIRES in the preamble (all already present in main.tex):
%%   amsmath, amssymb, booktabs, siunitx, multirow, array
%% =====================================================================

% ---- Notation macros (safe to move to the preamble) -----------------
\providecommand{\dJ}{\Delta J}
\providecommand{\Jt}{J_t}                  % composite inertia
\providecommand{\Jb}{J_b}                  % bare-airframe inertia
\providecommand{\cvec}{\mathbf{c}}
\providecommand{\rp}{\mathbf{r}_p}
\providecommand{\rc}{\mathbf{r}_c}
\providecommand{\mpay}{\hat{m}_p}          % payload-mass scalar (prior)
\providecommand{\mb}{m_b}                  % bare-airframe mass
\providecommand{\mt}{m_t}                  % composite mass
\providecommand{\Tphys}{T_{\mathrm{phys}}}
\providecommand{\tauphys}{\boldsymbol{\tau}_{\mathrm{phys}}}
\providecommand{\tcom}{\boldsymbol{\tau}_{\mathrm{com}}}
\providecommand{\ev}{\mathbf{e}_3}

\section{Modeling and Parameter Structure}\label{sec:model}

This section states the model the estimator and the controller share, and
then makes explicit \emph{which} of its parameters are estimated online.
The distinction is the substance of the method: a payload changes three
quantities in the dynamics --- mass, roll/pitch inertia, and the
horizontal offset of the composite centre of mass (CoM) --- yet only one
of them is a decision variable of the estimator. The remaining two are
generated from it and from known contact geometry by rigid-body algebra.

\subsection{Frames, notation, and conventions}\label{sec:model:frames}

We use an ENU world frame $\mathcal{W}$ and an FLU body frame
$\mathcal{B}$ whose origin is the geometric centre of the rotor disc.
The state and input are
\begin{equation}
x = \bigl[\mathbf{p};\, \mathbf{v};\, \mathbf{q};\, \boldsymbol{\omega}\bigr] \in \mathbb{R}^{13},
\qquad
u = \bigl[T;\, \boldsymbol{\tau}\bigr] \in \mathbb{R}^{4},
\label{eq:state}
\end{equation}
with position $\mathbf{p}$ and velocity $\mathbf{v}$ in $\mathcal{W}$,
body rate $\boldsymbol{\omega}$ in $\mathcal{B}$, collective thrust $T$
along the body $z$ axis, and body torque $\boldsymbol{\tau}$. We write
$\ev = [0,0,1]^{\!\top}$.

Attitude is a unit quaternion under the \emph{Hamilton} convention,
scalar first, $\mathbf{q} = [q_w, q_x, q_y, q_z]^{\!\top}$, representing
the rotation from $\mathcal{B}$ to $\mathcal{W}$; $R(\mathbf{q}) \in SO(3)$
denotes the corresponding rotation matrix. Stating the convention matters
because the body-rate kinematics \eqref{eq:quat} is a \emph{right}
multiplication precisely when $\boldsymbol{\omega}$ is resolved in
$\mathcal{B}$ --- the JPL convention, or a world-frame rate, would place
the product on the other side and integrate the attitude the wrong way.

\subsection{Rigid-body kinematics}\label{sec:model:kin}

Translational kinematics is trivial,
\begin{equation}
\dot{\mathbf{p}} = \mathbf{v}. \label{eq:kin}
\end{equation}
Attitude propagates by the quaternion product with the pure quaternion
$[0, \boldsymbol{\omega}]$,
\begin{equation}
\dot{\mathbf{q}} = \tfrac{1}{2}\, \mathbf{q} \otimes [0,\, \boldsymbol{\omega}]
= \tfrac{1}{2}
\begin{bmatrix}
-q_x & -q_y & -q_z \\
 q_w & -q_z &  q_y \\
 q_z &  q_w & -q_x \\
-q_y &  q_x &  q_w
\end{bmatrix} \boldsymbol{\omega},
\label{eq:quat}
\end{equation}
the two forms being algebraically identical; the implementation uses the
$4 \times 3$ matrix on the right. Neither \eqref{eq:kin} nor
\eqref{eq:quat} contains a payload-dependent parameter: \emph{the payload
enters the model only through the dynamics}, which is what confines the
parameter discussion below to three equations.

\subsection{Rigid-body dynamics}\label{sec:model:dyn}

\paragraph{Translation} Newton's law in $\mathcal{W}$, with a linear
aerodynamic drag term, reads
\begin{equation}
\dot{\mathbf{v}} = \frac{1}{m}\Bigl( R(\mathbf{q})\, \ev\, T - k_d\, \mathbf{v} \Bigr) - g\, \ev,
\label{eq:trans}
\end{equation}
where $m$ is the \emph{composite} mass (airframe plus payload) and $k_d$ a
fixed drag coefficient. The mass appears in exactly one role, as the
thrust-to-acceleration gain $1/m$. This is the mechanism of its
observability: given a thrust value reconstructed independently of the
model (\S\ref{sec:model:phys}), the measured $\dot{\mathbf{v}}$ determines
$m$ uniquely, and it does so in near-hover flight, without a dedicated
excitation manoeuvre.

\paragraph{Rotation} Let $\Jt$ denote the composite inertia, i.e.\ the
bare-airframe value $\Jb = \mathrm{diag}(J_{xx}, J_{yy}, J_{zz})$ raised by
the payload increment $\dJ$ on the roll and pitch axes,
\begin{equation}
\Jt = \mathrm{diag}\bigl(J_{xx} + \dJ,\; J_{yy} + \dJ,\; J_{zz}\bigr),
\label{eq:Jt}
\end{equation}
in the same spirit as the composite mass $\mt$ of
\S\ref{sec:model:payload}. Euler's equation about the composite CoM is
then
\begin{equation}
\dot{\boldsymbol{\omega}} = \Jt^{-1}\Bigl( \boldsymbol{\tau} + \tcom(\cvec) - \boldsymbol{\omega} \times \Jt \boldsymbol{\omega} \Bigr).
\label{eq:rot}
\end{equation}

\paragraph{Thrust eccentricity} The term $\tcom$ is the moment the
collective thrust exerts about the composite CoM. Thrust acts at the
rotor-disc centre --- the body origin --- along body $z$, while the CoM
sits at $\rc = [c_x, c_y, c_z]^{\!\top}$, so
\begin{equation}
\tcom(\cvec) = (-\rc) \times \ev T = \bigl[-c_y T,\; +c_x T,\; 0\bigr]^{\!\top}.
\label{eq:tcom}
\end{equation}
Two consequences are worth stating explicitly. First, $c_z$ \emph{cannot}
appear: a vertical CoM offset is collinear with the body-$z$ thrust, so
their cross product vanishes identically, and gravity acts at the CoM
itself and exerts no moment about it. The CoM offset is a two-parameter
quantity \emph{by construction, not by truncation} --- a point worth making
because ``the offset is three-dimensional yet only two components are
estimated'' is a natural objection. Second, \eqref{eq:tcom} is a
\emph{constant} moment in level hover, and a large one: at the payload of
\S\ref{sec:model:payload} it consumes roughly two thirds of the available
roll authority. Hovering on the spot with an off-centre load therefore
\emph{requires} a persistent trim torque, a fact the cost function must
also admit (\S\ref{sec:model:ctrl}).

Equations \eqref{eq:trans}--\eqref{eq:tcom} exhibit the property this
paper builds on: \textbf{each payload-dependent parameter has exactly one
point of entry.} The mass enters only through $1/m$ in \eqref{eq:trans},
the inertia increment only through $\Jt$ in \eqref{eq:rot}, and the CoM
offset only through \eqref{eq:tcom}.

\subsection{Payload-induced parameters}\label{sec:model:payload}

We model the payload as a point mass $\mpay$ rigidly attached at a known
body-frame offset $\rp = [r_x, r_y, r_z]^{\!\top}$ (with $r_z < 0$), and
write the composite mass and the reduced mass as
\begin{equation}
\mt = \mb + \mpay, \qquad \mu = \frac{\mb\, \mpay}{\mt}.
\label{eq:mu}
\end{equation}

\paragraph{CoM offset} The composite CoM is the mass-weighted mean
position, $\rc = (\mpay/\mt)\, \rp$, whence
\begin{equation}
\cvec = [c_x,\, c_y] = \frac{\mpay}{\mt}\, [r_x,\, r_y].
\label{eq:cxy}
\end{equation}

\paragraph{Inertia increment} Because the CoM \emph{moves} when the
payload is attached, the parallel-axis theorem must be applied about the
new CoM, and \emph{both} bodies contribute: the airframe now sits at
distance $(\mpay/\mt)\|\rp\|$ from it and the payload at
$(\mb/\mt)\|\rp\|$. For a lever arm $d$ the two terms collapse onto the
reduced mass,
\begin{equation}
\mb \Bigl(\frac{\mpay}{\mt}\Bigr)^{\!2} d^2 + \mpay \Bigl(\frac{\mb}{\mt}\Bigr)^{\!2} d^2
= \frac{\mb \mpay}{\mt}\, d^2 = \mu\, d^2 ,
\label{eq:reduced}
\end{equation}
which is why $\mu$ and not $\mpay$ is the correct coefficient. Writing
$\mpay d^2$ instead --- the natural slip, treating the payload as orbiting
a fixed body origin --- overestimates the increment by the factor
$\mt/\mb$, here \SI{14.5}{\percent}. Applying \eqref{eq:reduced} per axis
and taking the roll/pitch mean gives the single scalar the model carries,
\begin{align}
\dJ_{xx} &= \mu\bigl(r_y^2 + r_z^2\bigr), \quad
\dJ_{yy} = \mu\bigl(r_x^2 + r_z^2\bigr), \nonumber\\
\dJ &= \tfrac{1}{2}\bigl(\dJ_{xx} + \dJ_{yy}\bigr) = \mu\Bigl( r_z^2 + \tfrac{1}{2}\bigl(r_x^2 + r_y^2\bigr) \Bigr).
\label{eq:dj}
\end{align}

\paragraph{One unknown, three numbers} Equations
\eqref{eq:cxy}--\eqref{eq:dj} produce all three payload-dependent
quantities from a \emph{single} unknown scalar $\mpay$ together with known
geometry $\rp$. The rigid-body constraint thus replaces two degrees of
freedom that a naive augmentation would have to identify from data; the
estimator is not asked to rediscover the parallel-axis theorem from noise.
Table~\ref{tab:params} collects the resulting identities.

\paragraph{Magnitudes} For a \SI{0.3}{\kg} payload at
$\rp = [0,\, 0.109,\, -0.47]\,$\si{\meter} on our platform
($\mb = \SI{2.0643}{\kg}$, $J_{xx} = \SI{0.0142}{\kg\meter\squared}$),
\eqref{eq:mu}--\eqref{eq:dj} give $\mu = \SI{0.262}{\kg}$ and
$\dJ = \SI{0.0594}{\kg\meter\squared}$, an inertia ratio
$(J_{xx}+\dJ)/J_{xx} = 5.18$, together with $c_y = \SI{13.8}{\milli\meter}$
and hence a constant hover trim torque $|c_y|\mt g = \SI{0.32}{\newton\meter}$,
about \SI{64}{\percent} of the roll authority. A \SI{0.3}{\kg} parcel is
therefore not a \SI{15}{\percent} perturbation of the mass: it is a
five-fold change of inertia plus a standing demand on two thirds of the
torque margin. \emph{Adapting the mass alone cannot account for either.}

\begin{table*}[t]
\caption{Model parameters: which are estimated online, which are algebraically
derived, and which are fixed. \emph{Only the mass is a decision variable of the
estimator}; the inertia increment and CoM offset enter the MHE as \emph{known}
runtime parameters, on the same footing as the measured thrust and torque.}
\label{tab:params}
\centering
\footnotesize
\begin{tabular}{@{}l p{0.17\textwidth} c p{0.25\textwidth} p{0.28\textwidth}@{}}
\toprule
Symbol & Meaning & Dim. & Identity & Enters the model through \\
\midrule
$m$ & composite mass (airframe + payload) & 1 & \textbf{MHE decision variable} ($\dot m = 0$, state 14) & translational dynamics \eqref{eq:trans}; cost hover baseline \eqref{eq:hover} \\
$\dJ$ & roll/pitch inertia increment & 1 & algebraically derived, \eqref{eq:dj} & composite inertia \eqref{eq:Jt} in \eqref{eq:rot}; gain scaling \eqref{eq:gain} \\
$\cvec = [c_x, c_y]$ & composite-CoM horizontal offset & 2 & derived \eqref{eq:cxy}, or algebraically inverted online \eqref{eq:conline} & thrust eccentricity torque \eqref{eq:tcom}; cost torque baseline \eqref{eq:hover} \\
\midrule
$\mpay$ & payload-mass scalar (\emph{the single unknown behind $\dJ$ and $\cvec$}) & 1 & weak operational prior (nominal parcel mass) & \eqref{eq:cxy}, \eqref{eq:dj} \\
$\rp = [r_x, r_y, r_z]$ & payload offset in body frame & 3 & known (gripper contact geometry) & \eqref{eq:cxy}, \eqref{eq:dj} \\
\midrule
$\mb$ & bare-airframe mass & 1 & fixed calibration & \SI{2.0643}{\kg} \\
$\Jb$ & bare-airframe inertia & 3 & fixed calibration & $\mathrm{diag}(0.0142, 0.0142, 0.0210)$ \si{\kg\meter\squared} \\
$k_d$, $g$ & drag coefficient, gravity & 2 & fixed constants & $0.05$, \SI{9.81}{\meter\per\second\squared} \\
\bottomrule
\end{tabular}
\end{table*}

\subsection{Model-independent reconstruction of the physical input}\label{sec:model:phys}

Both $T$ and $\boldsymbol{\tau}$ in \eqref{eq:trans}--\eqref{eq:rot} are
reconstructed from measured rotor speeds $\omega_i$ and the known rotor
positions $(x_i, y_i)$, with $F_i = k_f \omega_i^2$,
\begin{equation}
\Tphys = \sum_i F_i, \quad
\tau_{\mathrm{roll}} = \sum_i y_i F_i, \quad
\tau_{\mathrm{pitch}} = -\sum_i x_i F_i .
\label{eq:phys}
\end{equation}
Feeding the controller's \emph{commanded} thrust back into the estimator
instead would create a self-consistent blind spot: at hover
$m \approx T/g$ holds for \emph{any} $m$ if $T$ is itself computed from the
assumed $m$, and the mass becomes unobservable. Equation \eqref{eq:phys}
breaks the loop because it reads a physical signal that no model state
enters.

\subsection{Estimation: the mass as the single online degree of freedom}\label{sec:model:mhe}

The moving-horizon estimator augments the mass into the state,
\begin{equation}
x_{\mathrm{aug}} = [x;\, m] \in \mathbb{R}^{14}, \qquad \dot m = 0,
\label{eq:xaug}
\end{equation}
holding it constant over the window so that the horizon of measurements
constrains one value rather than a free trajectory. Process noise
$\mathbf{w} \in \mathbb{R}^{13}$ acts on the physical states only; adding
noise on the mass channel would let it drift within the window and dissolve
the identification constraint. Over a window of $N$ stages the estimator
solves
\begin{align}
\min_{x_{0}, \mathbf{w}} \quad
& \sum_{k=0}^{N-1} \Bigl( \| \mathbf{y}_k - h(x_k) \|^2_{R} + \| \mathbf{w}_k \|^2_{Q} \Bigr)
+ \| x_0 - \bar{x}_0 \|^2_{Q_0} \nonumber \\
\text{s.t.} \quad & x_{k+1} = f\bigl(x_k, \mathbf{w}_k;\, p_{\mathrm{MHE}}\bigr), \quad
m \in [\,\underline{m},\, \overline{m}\,],
\label{eq:mhe}
\end{align}
with $f$ the discretisation of \eqref{eq:kin}--\eqref{eq:rot}, arrival cost
$\| \cdot \|^2_{Q_0}$ anchoring the window start, and measurement
$\mathbf{y}$ the 13 physical states. Crucially, the payload geometry is
\emph{not} in the decision vector but in the runtime parameter vector,
alongside the measured input:
\begin{equation}
p_{\mathrm{MHE}} = \bigl[\, T,\; \boldsymbol{\tau},\; \dJ,\; c_x,\; c_y \,\bigr] \in \mathbb{R}^{7}.
\label{eq:pmhe}
\end{equation}
Inside the estimator, $\dJ$ and $\cvec$ are \emph{known inputs}, on the
same footing as $T$ and $\boldsymbol{\tau}$. We use $N = 20$ at
$dt = \SI{0.1}{\second}$ (a \SI{2.0}{\second} horizon) and bound the mass
to $[\SI{1.961}{\kg}, \SI{5.0}{\kg}]$; the lower bound is the physical
floor $0.95\,\mb$, since the vehicle cannot weigh less than its own
airframe. The bound is a hard constraint on the decision variable --- an
advantage of the moving-horizon formulation over a recursive filter, which
can only be nudged away from non-physical estimates.

\paragraph{Ground-truth-free CoM inversion} When rotor speeds are
available the CoM offset need not come from the prior at all. In steady
hover the eccentric moment \eqref{eq:tcom} is balanced by the inner loop,
so $\dot{\boldsymbol{\omega}} \approx 0$ and \eqref{eq:rot} inverts
directly:
\begin{equation}
c_x = -\,\tau_{\mathrm{pitch}}^{\mathrm{phys}} / \Tphys, \qquad
c_y = +\,\tau_{\mathrm{roll}}^{\mathrm{phys}} / \Tphys .
\label{eq:conline}
\end{equation}
This is an algebraic inversion, not an optimisation, and it adds no degree
of freedom. Three properties make it admissible. \emph{(i)} It reads no
model state, so it cannot close an estimate-to-model feedback loop.
\emph{(ii)} The rotor coefficient $k_f$ cancels in the ratio ---
$c_x = \sum_i x_i \omega_i^2 / \sum_i \omega_i^2$ --- so the result is
immune to thrust-calibration error, unlike any scheme that consumes an
absolute thrust. \emph{(iii)} Channels are assigned by identifiability:
$\cvec$ has an independent observable and is estimated online, whereas
$r_z$ is annihilated by the cross product in \eqref{eq:tcom} and cannot be
identified from the torque channel at all --- it must remain a prior.
Samples are gated to steady flight
($\|\boldsymbol{\omega}\| < \SI{0.15}{\radian\per\second}$,
$\|\mathbf{v}_{xy}\| < \SI{0.20}{\meter\per\second}$), which rejects the
\emph{systematic} error of manoeuvring torques, and then low-pass filtered
with an exponential moving average of time constant \SI{2}{\second}, which
suppresses \emph{random} error: $\cvec$ is quasi-static once the load is
grasped, so its signal lies at DC while the differential-mode rotor noise
does not.

\subsection{Control: how the parameters re-centre the cost}\label{sec:model:ctrl}

The NMPC carries the parameters as runtime values,
\begin{equation}
p_{\mathrm{NMPC}} = \bigl[\, x_{\mathrm{ref}}(13),\; m,\; \dJ,\; \cvec(2),\; \mathbf{d}(3) \,\bigr] \in \mathbb{R}^{20},
\label{eq:pnmpc}
\end{equation}
where $\mathbf{d}$ is a lumped translational disturbance used only by the
$\mathcal{L}_1$ baseline of \S\ref{sec:results}-E and identically zero
under the proposed method. It minimises, over a horizon of $N$ stages,
\begin{equation}
\sum_{k} \Bigl( \| e_{\mathrm{track}}(x_k, x_{\mathrm{ref},k}) \|^2_{Q}
+ \| u_k - u_{\mathrm{hover}} \|^2_{R} \Bigr)
\label{eq:nmpc}
\end{equation}
subject to \eqref{eq:kin}--\eqref{eq:rot} and to input bounds on thrust and
torque. The second term penalises the deviation of the input from a
\emph{baseline}, and that baseline is recomputed from the live parameters
rather than fixed at build time:
\begin{equation}
T_h = m\,g, \qquad
u_{\mathrm{hover}} = \bigl[\, T_h,\; c_y T_h,\; -c_x T_h,\; 0 \,\bigr]^{\!\top}.
\label{eq:hover}
\end{equation}
Both entries matter, and for the same reason. Penalising $\|u\|^2_R$
outright would treat the thrust needed to hover as an expense to be
minimised; anchoring instead at a \emph{stale} $T_h$ merely moves the
error --- with a baseline computed from the unloaded mass, the cost pulls
the thrust down persistently and the tracking error settles at a nonzero
offset that no amount of retuning removes (increasing $R$ makes it worse,
since $R$ is the strength with which the wrong baseline pulls). The torque
entries are exactly $-\tcom$ of \eqref{eq:tcom}: hovering with an eccentric
load requires a constant trim torque, so a cost that does not admit it
pulls the torque toward zero and produces the same class of steady-state
bias. Equation \eqref{eq:tcom} teaches the \emph{model} that the moment
exists; equation \eqref{eq:hover} teaches the \emph{objective} that it is
necessary. Both are required.

\paragraph{Inner-loop bandwidth} Model consistency in the outer loop
cannot restore inner-loop bandwidth. A payload that multiplies roll/pitch
inertia lowers the effective bandwidth of the PX4 rate controller, and for
heavy eccentric loads the resulting attitude oscillation grows into torque
saturation. We therefore scale the rate gains at attach by the inertia
ratio,
\begin{equation}
K \leftarrow K_{\mathrm{nom}} \cdot \min\!\Bigl( \frac{J_{xx} + \dJ}{J_{xx}},\; 5.0 \Bigr).
\label{eq:gain}
\end{equation}
The cap is not incidental: at \SI{0.3}{\kg} the ratio is already $5.18$,
so that payload sits at the ceiling of what inner-loop retuning alone can
absorb on this airframe --- a boundary we quantify in
\S\ref{sec:results}-C.

\subsection{Why the inertia and CoM are not augmented into the state}\label{sec:model:whynot}

Augmenting $\dJ$ and $\cvec$ into the MHE decision vector is the obvious
alternative. We argue against it on four grounds.

\emph{(i) Asymmetric observability.} The mass is excited directly by the
thrust--acceleration relation \eqref{eq:trans} and is observable in
near-hover flight. The inertia increment acts only through
$\dot{\boldsymbol{\omega}}$ in \eqref{eq:rot}, which a delivery task ---
spent largely in near-hover, low-rate flight --- excites weakly. Making it
free lets the optimiser explain measurement noise with a nearly
unobservable quantity.

\emph{(ii) They compete for one residual.} Deriving the geometry from the
\emph{instantaneous} mass estimate closes a feedback loop
$m_{\mathrm{est}} \rightarrow \dJ, \cvec \rightarrow \text{attitude model}
\rightarrow m_{\mathrm{est}}$. We observed this bootstrap fail structurally
on light payloads: before lift-off the mass estimate reads low
\emph{legitimately}, that value is consumed as ``no geometry'', the
geometry channel closes, and after lift-off the model is mismatched and
the estimate sticks at its lower bound (\S\ref{sec:method}-E). The fix was
to sever the loop, not to add more free parameters.

\emph{(iii) Economy of parameterisation.} Equations
\eqref{eq:cxy}--\eqref{eq:dj} obtain three numbers from one scalar. Two
degrees of freedom are bought back by a constraint that holds exactly for a
rigid attachment.

\emph{(iv) The remaining prior may be weak.} With the geometry prior set
\SI{33}{\percent} wrong, the mass estimate still converges to within
\SI{0.3}{\percent} (\S\ref{sec:results}-C), so a nominal parcel mass
suffices on hardware and no attachment ground truth is required. Where
rotor speeds are available, \eqref{eq:conline} removes the prior from the
$\cvec$ path entirely.

\subsection{Declared modeling approximations}\label{sec:model:approx}

Three points deserve explicit statement rather than silence.

\emph{(a) Anisotropy of $\dJ$.} The single scalar in \eqref{eq:dj}
replaces the pair $(\dJ_{xx}, \dJ_{yy})$ by their mean, an error of
$\tfrac{1}{2}\mu|r_x^2 - r_y^2|$, below \SI{1}{\percent} of $J_{xx}+\dJ$ at
our offsets, and exactly zero when $r_x = r_y$.

\emph{(b) The yaw increment is neglected.} $\dJ_{zz} = \mu(r_x^2 + r_y^2)$
is not added to $J_{zz}$ in \eqref{eq:Jt}. At \SI{0.3}{\kg} and
$r_{xy} = \SI{0.109}{\meter}$ it amounts to \SI{0.0031}{\kg\meter\squared},
some \SI{15}{\percent} of $J_{zz}$ --- a known, non-negligible
approximation. It does not affect the reported results, since the yaw
channel neither carries payload compensation nor participates in
\eqref{eq:conline}; and removing it requires only substituting
$J_{zz} + \mu(r_x^2 + r_y^2)$ in \eqref{eq:Jt}, introducing no new
estimated quantity.

\emph{(c) $c_z$ is absent by construction.} As shown below
\eqref{eq:tcom}, this is an exact algebraic consequence, not a truncation,
and is listed here only to forestall the opposite reading.
   (place after \section{Related Work};
%%         it supersedes the "Dynamics and parameterization" / "MHE" /
%%         "Inner-loop gain scaling" paragraphs of the current
%%         \section{System Overview and Problem Setup}, and subsumes
%%         sec3_parameters.tex in full.)
%%
%% REQUIRES in the preamble (all already present in main.tex):
%%   amsmath, amssymb, booktabs, siunitx, multirow, array
%% =====================================================================

% ---- Notation macros (safe to move to the preamble) -----------------
\providecommand{\dJ}{\Delta J}
\providecommand{\Jt}{J_t}                  % composite inertia
\providecommand{\Jb}{J_b}                  % bare-airframe inertia
\providecommand{\cvec}{\mathbf{c}}
\providecommand{\rp}{\mathbf{r}_p}
\providecommand{\rc}{\mathbf{r}_c}
\providecommand{\mpay}{\hat{m}_p}          % payload-mass scalar (prior)
\providecommand{\mb}{m_b}                  % bare-airframe mass
\providecommand{\mt}{m_t}                  % composite mass
\providecommand{\Tphys}{T_{\mathrm{phys}}}
\providecommand{\tauphys}{\boldsymbol{\tau}_{\mathrm{phys}}}
\providecommand{\tcom}{\boldsymbol{\tau}_{\mathrm{com}}}
\providecommand{\ev}{\mathbf{e}_3}

\section{Modeling and Parameter Structure}\label{sec:model}

This section states the model the estimator and the controller share, and
then makes explicit \emph{which} of its parameters are estimated online.
The distinction is the substance of the method: a payload changes three
quantities in the dynamics --- mass, roll/pitch inertia, and the
horizontal offset of the composite centre of mass (CoM) --- yet only one
of them is a decision variable of the estimator. The remaining two are
generated from it and from known contact geometry by rigid-body algebra.

\subsection{Frames, notation, and conventions}\label{sec:model:frames}

We use an ENU world frame $\mathcal{W}$ and an FLU body frame
$\mathcal{B}$ whose origin is the geometric centre of the rotor disc.
The state and input are
\begin{equation}
x = \bigl[\mathbf{p};\, \mathbf{v};\, \mathbf{q};\, \boldsymbol{\omega}\bigr] \in \mathbb{R}^{13},
\qquad
u = \bigl[T;\, \boldsymbol{\tau}\bigr] \in \mathbb{R}^{4},
\label{eq:state}
\end{equation}
with position $\mathbf{p}$ and velocity $\mathbf{v}$ in $\mathcal{W}$,
body rate $\boldsymbol{\omega}$ in $\mathcal{B}$, collective thrust $T$
along the body $z$ axis, and body torque $\boldsymbol{\tau}$. We write
$\ev = [0,0,1]^{\!\top}$.

Attitude is a unit quaternion under the \emph{Hamilton} convention,
scalar first, $\mathbf{q} = [q_w, q_x, q_y, q_z]^{\!\top}$, representing
the rotation from $\mathcal{B}$ to $\mathcal{W}$; $R(\mathbf{q}) \in SO(3)$
denotes the corresponding rotation matrix. Stating the convention matters
because the body-rate kinematics \eqref{eq:quat} is a \emph{right}
multiplication precisely when $\boldsymbol{\omega}$ is resolved in
$\mathcal{B}$ --- the JPL convention, or a world-frame rate, would place
the product on the other side and integrate the attitude the wrong way.

\subsection{Rigid-body kinematics}\label{sec:model:kin}

Translational kinematics is trivial,
\begin{equation}
\dot{\mathbf{p}} = \mathbf{v}. \label{eq:kin}
\end{equation}
Attitude propagates by the quaternion product with the pure quaternion
$[0, \boldsymbol{\omega}]$,
\begin{equation}
\dot{\mathbf{q}} = \tfrac{1}{2}\, \mathbf{q} \otimes [0,\, \boldsymbol{\omega}]
= \tfrac{1}{2}
\begin{bmatrix}
-q_x & -q_y & -q_z \\
 q_w & -q_z &  q_y \\
 q_z &  q_w & -q_x \\
-q_y &  q_x &  q_w
\end{bmatrix} \boldsymbol{\omega},
\label{eq:quat}
\end{equation}
the two forms being algebraically identical; the implementation uses the
$4 \times 3$ matrix on the right. Neither \eqref{eq:kin} nor
\eqref{eq:quat} contains a payload-dependent parameter: \emph{the payload
enters the model only through the dynamics}, which is what confines the
parameter discussion below to three equations.

\subsection{Rigid-body dynamics}\label{sec:model:dyn}

\paragraph{Translation} Newton's law in $\mathcal{W}$, with a linear
aerodynamic drag term, reads
\begin{equation}
\dot{\mathbf{v}} = \frac{1}{m}\Bigl( R(\mathbf{q})\, \ev\, T - k_d\, \mathbf{v} \Bigr) - g\, \ev,
\label{eq:trans}
\end{equation}
where $m$ is the \emph{composite} mass (airframe plus payload) and $k_d$ a
fixed drag coefficient. The mass appears in exactly one role, as the
thrust-to-acceleration gain $1/m$. This is the mechanism of its
observability: given a thrust value reconstructed independently of the
model (\S\ref{sec:model:phys}), the measured $\dot{\mathbf{v}}$ determines
$m$ uniquely, and it does so in near-hover flight, without a dedicated
excitation manoeuvre.

\paragraph{Rotation} Let $\Jt$ denote the composite inertia, i.e.\ the
bare-airframe value $\Jb = \mathrm{diag}(J_{xx}, J_{yy}, J_{zz})$ raised by
the payload increment $\dJ$ on the roll and pitch axes,
\begin{equation}
\Jt = \mathrm{diag}\bigl(J_{xx} + \dJ,\; J_{yy} + \dJ,\; J_{zz}\bigr),
\label{eq:Jt}
\end{equation}
in the same spirit as the composite mass $\mt$ of
\S\ref{sec:model:payload}. Euler's equation about the composite CoM is
then
\begin{equation}
\dot{\boldsymbol{\omega}} = \Jt^{-1}\Bigl( \boldsymbol{\tau} + \tcom(\cvec) - \boldsymbol{\omega} \times \Jt \boldsymbol{\omega} \Bigr).
\label{eq:rot}
\end{equation}

\paragraph{Thrust eccentricity} The term $\tcom$ is the moment the
collective thrust exerts about the composite CoM. Thrust acts at the
rotor-disc centre --- the body origin --- along body $z$, while the CoM
sits at $\rc = [c_x, c_y, c_z]^{\!\top}$, so
\begin{equation}
\tcom(\cvec) = (-\rc) \times \ev T = \bigl[-c_y T,\; +c_x T,\; 0\bigr]^{\!\top}.
\label{eq:tcom}
\end{equation}
Two consequences are worth stating explicitly. First, $c_z$ \emph{cannot}
appear: a vertical CoM offset is collinear with the body-$z$ thrust, so
their cross product vanishes identically, and gravity acts at the CoM
itself and exerts no moment about it. The CoM offset is a two-parameter
quantity \emph{by construction, not by truncation} --- a point worth making
because ``the offset is three-dimensional yet only two components are
estimated'' is a natural objection. Second, \eqref{eq:tcom} is a
\emph{constant} moment in level hover, and a large one: at the payload of
\S\ref{sec:model:payload} it consumes roughly two thirds of the available
roll authority. Hovering on the spot with an off-centre load therefore
\emph{requires} a persistent trim torque, a fact the cost function must
also admit (\S\ref{sec:model:ctrl}).

Equations \eqref{eq:trans}--\eqref{eq:tcom} exhibit the property this
paper builds on: \textbf{each payload-dependent parameter has exactly one
point of entry.} The mass enters only through $1/m$ in \eqref{eq:trans},
the inertia increment only through $\Jt$ in \eqref{eq:rot}, and the CoM
offset only through \eqref{eq:tcom}.

\subsection{Payload-induced parameters}\label{sec:model:payload}

We model the payload as a point mass $\mpay$ rigidly attached at a known
body-frame offset $\rp = [r_x, r_y, r_z]^{\!\top}$ (with $r_z < 0$), and
write the composite mass and the reduced mass as
\begin{equation}
\mt = \mb + \mpay, \qquad \mu = \frac{\mb\, \mpay}{\mt}.
\label{eq:mu}
\end{equation}

\paragraph{CoM offset} The composite CoM is the mass-weighted mean
position, $\rc = (\mpay/\mt)\, \rp$, whence
\begin{equation}
\cvec = [c_x,\, c_y] = \frac{\mpay}{\mt}\, [r_x,\, r_y].
\label{eq:cxy}
\end{equation}

\paragraph{Inertia increment} Because the CoM \emph{moves} when the
payload is attached, the parallel-axis theorem must be applied about the
new CoM, and \emph{both} bodies contribute: the airframe now sits at
distance $(\mpay/\mt)\|\rp\|$ from it and the payload at
$(\mb/\mt)\|\rp\|$. For a lever arm $d$ the two terms collapse onto the
reduced mass,
\begin{equation}
\mb \Bigl(\frac{\mpay}{\mt}\Bigr)^{\!2} d^2 + \mpay \Bigl(\frac{\mb}{\mt}\Bigr)^{\!2} d^2
= \frac{\mb \mpay}{\mt}\, d^2 = \mu\, d^2 ,
\label{eq:reduced}
\end{equation}
which is why $\mu$ and not $\mpay$ is the correct coefficient. Writing
$\mpay d^2$ instead --- the natural slip, treating the payload as orbiting
a fixed body origin --- overestimates the increment by the factor
$\mt/\mb$, here \SI{14.5}{\percent}. Applying \eqref{eq:reduced} per axis
and taking the roll/pitch mean gives the single scalar the model carries,
\begin{align}
\dJ_{xx} &= \mu\bigl(r_y^2 + r_z^2\bigr), \quad
\dJ_{yy} = \mu\bigl(r_x^2 + r_z^2\bigr), \nonumber\\
\dJ &= \tfrac{1}{2}\bigl(\dJ_{xx} + \dJ_{yy}\bigr) = \mu\Bigl( r_z^2 + \tfrac{1}{2}\bigl(r_x^2 + r_y^2\bigr) \Bigr).
\label{eq:dj}
\end{align}

\paragraph{One unknown, three numbers} Equations
\eqref{eq:cxy}--\eqref{eq:dj} produce all three payload-dependent
quantities from a \emph{single} unknown scalar $\mpay$ together with known
geometry $\rp$. The rigid-body constraint thus replaces two degrees of
freedom that a naive augmentation would have to identify from data; the
estimator is not asked to rediscover the parallel-axis theorem from noise.
Table~\ref{tab:params} collects the resulting identities.

\paragraph{Magnitudes} For a \SI{0.3}{\kg} payload at
$\rp = [0,\, 0.109,\, -0.47]\,$\si{\meter} on our platform
($\mb = \SI{2.0643}{\kg}$, $J_{xx} = \SI{0.0142}{\kg\meter\squared}$),
\eqref{eq:mu}--\eqref{eq:dj} give $\mu = \SI{0.262}{\kg}$ and
$\dJ = \SI{0.0594}{\kg\meter\squared}$, an inertia ratio
$(J_{xx}+\dJ)/J_{xx} = 5.18$, together with $c_y = \SI{13.8}{\milli\meter}$
and hence a constant hover trim torque $|c_y|\mt g = \SI{0.32}{\newton\meter}$,
about \SI{64}{\percent} of the roll authority. A \SI{0.3}{\kg} parcel is
therefore not a \SI{15}{\percent} perturbation of the mass: it is a
five-fold change of inertia plus a standing demand on two thirds of the
torque margin. \emph{Adapting the mass alone cannot account for either.}

\begin{table*}[t]
\caption{Model parameters: which are estimated online, which are algebraically
derived, and which are fixed. \emph{Only the mass is a decision variable of the
estimator}; the inertia increment and CoM offset enter the MHE as \emph{known}
runtime parameters, on the same footing as the measured thrust and torque.}
\label{tab:params}
\centering
\footnotesize
\begin{tabular}{@{}l p{0.17\textwidth} c p{0.25\textwidth} p{0.28\textwidth}@{}}
\toprule
Symbol & Meaning & Dim. & Identity & Enters the model through \\
\midrule
$m$ & composite mass (airframe + payload) & 1 & \textbf{MHE decision variable} ($\dot m = 0$, state 14) & translational dynamics \eqref{eq:trans}; cost hover baseline \eqref{eq:hover} \\
$\dJ$ & roll/pitch inertia increment & 1 & algebraically derived, \eqref{eq:dj} & composite inertia \eqref{eq:Jt} in \eqref{eq:rot}; gain scaling \eqref{eq:gain} \\
$\cvec = [c_x, c_y]$ & composite-CoM horizontal offset & 2 & derived \eqref{eq:cxy}, or algebraically inverted online \eqref{eq:conline} & thrust eccentricity torque \eqref{eq:tcom}; cost torque baseline \eqref{eq:hover} \\
\midrule
$\mpay$ & payload-mass scalar (\emph{the single unknown behind $\dJ$ and $\cvec$}) & 1 & weak operational prior (nominal parcel mass) & \eqref{eq:cxy}, \eqref{eq:dj} \\
$\rp = [r_x, r_y, r_z]$ & payload offset in body frame & 3 & known (gripper contact geometry) & \eqref{eq:cxy}, \eqref{eq:dj} \\
\midrule
$\mb$ & bare-airframe mass & 1 & fixed calibration & \SI{2.0643}{\kg} \\
$\Jb$ & bare-airframe inertia & 3 & fixed calibration & $\mathrm{diag}(0.0142, 0.0142, 0.0210)$ \si{\kg\meter\squared} \\
$k_d$, $g$ & drag coefficient, gravity & 2 & fixed constants & $0.05$, \SI{9.81}{\meter\per\second\squared} \\
\bottomrule
\end{tabular}
\end{table*}

\subsection{Model-independent reconstruction of the physical input}\label{sec:model:phys}

Both $T$ and $\boldsymbol{\tau}$ in \eqref{eq:trans}--\eqref{eq:rot} are
reconstructed from measured rotor speeds $\omega_i$ and the known rotor
positions $(x_i, y_i)$, with $F_i = k_f \omega_i^2$,
\begin{equation}
\Tphys = \sum_i F_i, \quad
\tau_{\mathrm{roll}} = \sum_i y_i F_i, \quad
\tau_{\mathrm{pitch}} = -\sum_i x_i F_i .
\label{eq:phys}
\end{equation}
Feeding the controller's \emph{commanded} thrust back into the estimator
instead would create a self-consistent blind spot: at hover
$m \approx T/g$ holds for \emph{any} $m$ if $T$ is itself computed from the
assumed $m$, and the mass becomes unobservable. Equation \eqref{eq:phys}
breaks the loop because it reads a physical signal that no model state
enters.

\subsection{Estimation: the mass as the single online degree of freedom}\label{sec:model:mhe}

The moving-horizon estimator augments the mass into the state,
\begin{equation}
x_{\mathrm{aug}} = [x;\, m] \in \mathbb{R}^{14}, \qquad \dot m = 0,
\label{eq:xaug}
\end{equation}
holding it constant over the window so that the horizon of measurements
constrains one value rather than a free trajectory. Process noise
$\mathbf{w} \in \mathbb{R}^{13}$ acts on the physical states only; adding
noise on the mass channel would let it drift within the window and dissolve
the identification constraint. Over a window of $N$ stages the estimator
solves
\begin{align}
\min_{x_{0}, \mathbf{w}} \quad
& \sum_{k=0}^{N-1} \Bigl( \| \mathbf{y}_k - h(x_k) \|^2_{R} + \| \mathbf{w}_k \|^2_{Q} \Bigr)
+ \| x_0 - \bar{x}_0 \|^2_{Q_0} \nonumber \\
\text{s.t.} \quad & x_{k+1} = f\bigl(x_k, \mathbf{w}_k;\, p_{\mathrm{MHE}}\bigr), \quad
m \in [\,\underline{m},\, \overline{m}\,],
\label{eq:mhe}
\end{align}
with $f$ the discretisation of \eqref{eq:kin}--\eqref{eq:rot}, arrival cost
$\| \cdot \|^2_{Q_0}$ anchoring the window start, and measurement
$\mathbf{y}$ the 13 physical states. Crucially, the payload geometry is
\emph{not} in the decision vector but in the runtime parameter vector,
alongside the measured input:
\begin{equation}
p_{\mathrm{MHE}} = \bigl[\, T,\; \boldsymbol{\tau},\; \dJ,\; c_x,\; c_y \,\bigr] \in \mathbb{R}^{7}.
\label{eq:pmhe}
\end{equation}
Inside the estimator, $\dJ$ and $\cvec$ are \emph{known inputs}, on the
same footing as $T$ and $\boldsymbol{\tau}$. We use $N = 20$ at
$dt = \SI{0.1}{\second}$ (a \SI{2.0}{\second} horizon) and bound the mass
to $[\SI{1.961}{\kg}, \SI{5.0}{\kg}]$; the lower bound is the physical
floor $0.95\,\mb$, since the vehicle cannot weigh less than its own
airframe. The bound is a hard constraint on the decision variable --- an
advantage of the moving-horizon formulation over a recursive filter, which
can only be nudged away from non-physical estimates.

\paragraph{Ground-truth-free CoM inversion} When rotor speeds are
available the CoM offset need not come from the prior at all. In steady
hover the eccentric moment \eqref{eq:tcom} is balanced by the inner loop,
so $\dot{\boldsymbol{\omega}} \approx 0$ and \eqref{eq:rot} inverts
directly:
\begin{equation}
c_x = -\,\tau_{\mathrm{pitch}}^{\mathrm{phys}} / \Tphys, \qquad
c_y = +\,\tau_{\mathrm{roll}}^{\mathrm{phys}} / \Tphys .
\label{eq:conline}
\end{equation}
This is an algebraic inversion, not an optimisation, and it adds no degree
of freedom. Three properties make it admissible. \emph{(i)} It reads no
model state, so it cannot close an estimate-to-model feedback loop.
\emph{(ii)} The rotor coefficient $k_f$ cancels in the ratio ---
$c_x = \sum_i x_i \omega_i^2 / \sum_i \omega_i^2$ --- so the result is
immune to thrust-calibration error, unlike any scheme that consumes an
absolute thrust. \emph{(iii)} Channels are assigned by identifiability:
$\cvec$ has an independent observable and is estimated online, whereas
$r_z$ is annihilated by the cross product in \eqref{eq:tcom} and cannot be
identified from the torque channel at all --- it must remain a prior.
Samples are gated to steady flight
($\|\boldsymbol{\omega}\| < \SI{0.15}{\radian\per\second}$,
$\|\mathbf{v}_{xy}\| < \SI{0.20}{\meter\per\second}$), which rejects the
\emph{systematic} error of manoeuvring torques, and then low-pass filtered
with an exponential moving average of time constant \SI{2}{\second}, which
suppresses \emph{random} error: $\cvec$ is quasi-static once the load is
grasped, so its signal lies at DC while the differential-mode rotor noise
does not.

\subsection{Control: how the parameters re-centre the cost}\label{sec:model:ctrl}

The NMPC carries the parameters as runtime values,
\begin{equation}
p_{\mathrm{NMPC}} = \bigl[\, x_{\mathrm{ref}}(13),\; m,\; \dJ,\; \cvec(2),\; \mathbf{d}(3) \,\bigr] \in \mathbb{R}^{20},
\label{eq:pnmpc}
\end{equation}
where $\mathbf{d}$ is a lumped translational disturbance used only by the
$\mathcal{L}_1$ baseline of \S\ref{sec:results}-E and identically zero
under the proposed method. It minimises, over a horizon of $N$ stages,
\begin{equation}
\sum_{k} \Bigl( \| e_{\mathrm{track}}(x_k, x_{\mathrm{ref},k}) \|^2_{Q}
+ \| u_k - u_{\mathrm{hover}} \|^2_{R} \Bigr)
\label{eq:nmpc}
\end{equation}
subject to \eqref{eq:kin}--\eqref{eq:rot} and to input bounds on thrust and
torque. The second term penalises the deviation of the input from a
\emph{baseline}, and that baseline is recomputed from the live parameters
rather than fixed at build time:
\begin{equation}
T_h = m\,g, \qquad
u_{\mathrm{hover}} = \bigl[\, T_h,\; c_y T_h,\; -c_x T_h,\; 0 \,\bigr]^{\!\top}.
\label{eq:hover}
\end{equation}
Both entries matter, and for the same reason. Penalising $\|u\|^2_R$
outright would treat the thrust needed to hover as an expense to be
minimised; anchoring instead at a \emph{stale} $T_h$ merely moves the
error --- with a baseline computed from the unloaded mass, the cost pulls
the thrust down persistently and the tracking error settles at a nonzero
offset that no amount of retuning removes (increasing $R$ makes it worse,
since $R$ is the strength with which the wrong baseline pulls). The torque
entries are exactly $-\tcom$ of \eqref{eq:tcom}: hovering with an eccentric
load requires a constant trim torque, so a cost that does not admit it
pulls the torque toward zero and produces the same class of steady-state
bias. Equation \eqref{eq:tcom} teaches the \emph{model} that the moment
exists; equation \eqref{eq:hover} teaches the \emph{objective} that it is
necessary. Both are required.

\paragraph{Inner-loop bandwidth} Model consistency in the outer loop
cannot restore inner-loop bandwidth. A payload that multiplies roll/pitch
inertia lowers the effective bandwidth of the PX4 rate controller, and for
heavy eccentric loads the resulting attitude oscillation grows into torque
saturation. We therefore scale the rate gains at attach by the inertia
ratio,
\begin{equation}
K \leftarrow K_{\mathrm{nom}} \cdot \min\!\Bigl( \frac{J_{xx} + \dJ}{J_{xx}},\; 5.0 \Bigr).
\label{eq:gain}
\end{equation}
The cap is not incidental: at \SI{0.3}{\kg} the ratio is already $5.18$,
so that payload sits at the ceiling of what inner-loop retuning alone can
absorb on this airframe --- a boundary we quantify in
\S\ref{sec:results}-C.

\subsection{Why the inertia and CoM are not augmented into the state}\label{sec:model:whynot}

Augmenting $\dJ$ and $\cvec$ into the MHE decision vector is the obvious
alternative. We argue against it on four grounds.

\emph{(i) Asymmetric observability.} The mass is excited directly by the
thrust--acceleration relation \eqref{eq:trans} and is observable in
near-hover flight. The inertia increment acts only through
$\dot{\boldsymbol{\omega}}$ in \eqref{eq:rot}, which a delivery task ---
spent largely in near-hover, low-rate flight --- excites weakly. Making it
free lets the optimiser explain measurement noise with a nearly
unobservable quantity.

\emph{(ii) They compete for one residual.} Deriving the geometry from the
\emph{instantaneous} mass estimate closes a feedback loop
$m_{\mathrm{est}} \rightarrow \dJ, \cvec \rightarrow \text{attitude model}
\rightarrow m_{\mathrm{est}}$. We observed this bootstrap fail structurally
on light payloads: before lift-off the mass estimate reads low
\emph{legitimately}, that value is consumed as ``no geometry'', the
geometry channel closes, and after lift-off the model is mismatched and
the estimate sticks at its lower bound (\S\ref{sec:method}-E). The fix was
to sever the loop, not to add more free parameters.

\emph{(iii) Economy of parameterisation.} Equations
\eqref{eq:cxy}--\eqref{eq:dj} obtain three numbers from one scalar. Two
degrees of freedom are bought back by a constraint that holds exactly for a
rigid attachment.

\emph{(iv) The remaining prior may be weak.} With the geometry prior set
\SI{33}{\percent} wrong, the mass estimate still converges to within
\SI{0.3}{\percent} (\S\ref{sec:results}-C), so a nominal parcel mass
suffices on hardware and no attachment ground truth is required. Where
rotor speeds are available, \eqref{eq:conline} removes the prior from the
$\cvec$ path entirely.

\subsection{Declared modeling approximations}\label{sec:model:approx}

Three points deserve explicit statement rather than silence.

\emph{(a) Anisotropy of $\dJ$.} The single scalar in \eqref{eq:dj}
replaces the pair $(\dJ_{xx}, \dJ_{yy})$ by their mean, an error of
$\tfrac{1}{2}\mu|r_x^2 - r_y^2|$, below \SI{1}{\percent} of $J_{xx}+\dJ$ at
our offsets, and exactly zero when $r_x = r_y$.

\emph{(b) The yaw increment is neglected.} $\dJ_{zz} = \mu(r_x^2 + r_y^2)$
is not added to $J_{zz}$ in \eqref{eq:Jt}. At \SI{0.3}{\kg} and
$r_{xy} = \SI{0.109}{\meter}$ it amounts to \SI{0.0031}{\kg\meter\squared},
some \SI{15}{\percent} of $J_{zz}$ --- a known, non-negligible
approximation. It does not affect the reported results, since the yaw
channel neither carries payload compensation nor participates in
\eqref{eq:conline}; and removing it requires only substituting
$J_{zz} + \mu(r_x^2 + r_y^2)$ in \eqref{eq:Jt}, introducing no new
estimated quantity.

\emph{(c) $c_z$ is absent by construction.} As shown below
\eqref{eq:tcom}, this is an exact algebraic consequence, not a truncation,
and is listed here only to forestall the opposite reading.
   (place after \section{Related Work};
%%         it supersedes the "Dynamics and parameterization" / "MHE" /
%%         "Inner-loop gain scaling" paragraphs of the current
%%         \section{System Overview and Problem Setup}, and subsumes
%%         sec3_parameters.tex in full.)
%%
%% REQUIRES in the preamble (all already present in main.tex):
%%   amsmath, amssymb, booktabs, siunitx, multirow, array
%% =====================================================================

% ---- Notation macros (safe to move to the preamble) -----------------
\providecommand{\dJ}{\Delta J}
\providecommand{\Jt}{J_t}                  % composite inertia
\providecommand{\Jb}{J_b}                  % bare-airframe inertia
\providecommand{\cvec}{\mathbf{c}}
\providecommand{\rp}{\mathbf{r}_p}
\providecommand{\rc}{\mathbf{r}_c}
\providecommand{\mpay}{\hat{m}_p}          % payload-mass scalar (prior)
\providecommand{\mb}{m_b}                  % bare-airframe mass
\providecommand{\mt}{m_t}                  % composite mass
\providecommand{\Tphys}{T_{\mathrm{phys}}}
\providecommand{\tauphys}{\boldsymbol{\tau}_{\mathrm{phys}}}
\providecommand{\tcom}{\boldsymbol{\tau}_{\mathrm{com}}}
\providecommand{\ev}{\mathbf{e}_3}

\section{Modeling and Parameter Structure}\label{sec:model}

This section states the model the estimator and the controller share, and
then makes explicit \emph{which} of its parameters are estimated online.
The distinction is the substance of the method: a payload changes three
quantities in the dynamics --- mass, roll/pitch inertia, and the
horizontal offset of the composite centre of mass (CoM) --- yet only one
of them is a decision variable of the estimator. The remaining two are
generated from it and from known contact geometry by rigid-body algebra.

\subsection{Frames, notation, and conventions}\label{sec:model:frames}

We use an ENU world frame $\mathcal{W}$ and an FLU body frame
$\mathcal{B}$ whose origin is the geometric centre of the rotor disc.
The state and input are
\begin{equation}
x = \bigl[\mathbf{p};\, \mathbf{v};\, \mathbf{q};\, \boldsymbol{\omega}\bigr] \in \mathbb{R}^{13},
\qquad
u = \bigl[T;\, \boldsymbol{\tau}\bigr] \in \mathbb{R}^{4},
\label{eq:state}
\end{equation}
with position $\mathbf{p}$ and velocity $\mathbf{v}$ in $\mathcal{W}$,
body rate $\boldsymbol{\omega}$ in $\mathcal{B}$, collective thrust $T$
along the body $z$ axis, and body torque $\boldsymbol{\tau}$. We write
$\ev = [0,0,1]^{\!\top}$.

Attitude is a unit quaternion under the \emph{Hamilton} convention,
scalar first, $\mathbf{q} = [q_w, q_x, q_y, q_z]^{\!\top}$, representing
the rotation from $\mathcal{B}$ to $\mathcal{W}$; $R(\mathbf{q}) \in SO(3)$
denotes the corresponding rotation matrix. Stating the convention matters
because the body-rate kinematics \eqref{eq:quat} is a \emph{right}
multiplication precisely when $\boldsymbol{\omega}$ is resolved in
$\mathcal{B}$ --- the JPL convention, or a world-frame rate, would place
the product on the other side and integrate the attitude the wrong way.

\subsection{Rigid-body kinematics}\label{sec:model:kin}

Translational kinematics is trivial,
\begin{equation}
\dot{\mathbf{p}} = \mathbf{v}. \label{eq:kin}
\end{equation}
Attitude propagates by the quaternion product with the pure quaternion
$[0, \boldsymbol{\omega}]$,
\begin{equation}
\dot{\mathbf{q}} = \tfrac{1}{2}\, \mathbf{q} \otimes [0,\, \boldsymbol{\omega}]
= \tfrac{1}{2}
\begin{bmatrix}
-q_x & -q_y & -q_z \\
 q_w & -q_z &  q_y \\
 q_z &  q_w & -q_x \\
-q_y &  q_x &  q_w
\end{bmatrix} \boldsymbol{\omega},
\label{eq:quat}
\end{equation}
the two forms being algebraically identical; the implementation uses the
$4 \times 3$ matrix on the right. Neither \eqref{eq:kin} nor
\eqref{eq:quat} contains a payload-dependent parameter: \emph{the payload
enters the model only through the dynamics}, which is what confines the
parameter discussion below to three equations.

\subsection{Rigid-body dynamics}\label{sec:model:dyn}

\paragraph{Translation} Newton's law in $\mathcal{W}$, with a linear
aerodynamic drag term, reads
\begin{equation}
\dot{\mathbf{v}} = \frac{1}{m}\Bigl( R(\mathbf{q})\, \ev\, T - k_d\, \mathbf{v} \Bigr) - g\, \ev,
\label{eq:trans}
\end{equation}
where $m$ is the \emph{composite} mass (airframe plus payload) and $k_d$ a
fixed drag coefficient. The mass appears in exactly one role, as the
thrust-to-acceleration gain $1/m$. This is the mechanism of its
observability: given a thrust value reconstructed independently of the
model (\S\ref{sec:model:phys}), the measured $\dot{\mathbf{v}}$ determines
$m$ uniquely, and it does so in near-hover flight, without a dedicated
excitation manoeuvre.

\paragraph{Rotation} Let $\Jt$ denote the composite inertia, i.e.\ the
bare-airframe value $\Jb = \mathrm{diag}(J_{xx}, J_{yy}, J_{zz})$ raised by
the payload increment $\dJ$ on the roll and pitch axes,
\begin{equation}
\Jt = \mathrm{diag}\bigl(J_{xx} + \dJ,\; J_{yy} + \dJ,\; J_{zz}\bigr),
\label{eq:Jt}
\end{equation}
in the same spirit as the composite mass $\mt$ of
\S\ref{sec:model:payload}. Euler's equation about the composite CoM is
then
\begin{equation}
\dot{\boldsymbol{\omega}} = \Jt^{-1}\Bigl( \boldsymbol{\tau} + \tcom(\cvec) - \boldsymbol{\omega} \times \Jt \boldsymbol{\omega} \Bigr).
\label{eq:rot}
\end{equation}

\paragraph{Thrust eccentricity} The term $\tcom$ is the moment the
collective thrust exerts about the composite CoM. Thrust acts at the
rotor-disc centre --- the body origin --- along body $z$, while the CoM
sits at $\rc = [c_x, c_y, c_z]^{\!\top}$, so
\begin{equation}
\tcom(\cvec) = (-\rc) \times \ev T = \bigl[-c_y T,\; +c_x T,\; 0\bigr]^{\!\top}.
\label{eq:tcom}
\end{equation}
Two consequences are worth stating explicitly. First, $c_z$ \emph{cannot}
appear: a vertical CoM offset is collinear with the body-$z$ thrust, so
their cross product vanishes identically, and gravity acts at the CoM
itself and exerts no moment about it. The CoM offset is a two-parameter
quantity \emph{by construction, not by truncation} --- a point worth making
because ``the offset is three-dimensional yet only two components are
estimated'' is a natural objection. Second, \eqref{eq:tcom} is a
\emph{constant} moment in level hover, and a large one: at the payload of
\S\ref{sec:model:payload} it consumes roughly two thirds of the available
roll authority. Hovering on the spot with an off-centre load therefore
\emph{requires} a persistent trim torque, a fact the cost function must
also admit (\S\ref{sec:model:ctrl}).

Equations \eqref{eq:trans}--\eqref{eq:tcom} exhibit the property this
paper builds on: \textbf{each payload-dependent parameter has exactly one
point of entry.} The mass enters only through $1/m$ in \eqref{eq:trans},
the inertia increment only through $\Jt$ in \eqref{eq:rot}, and the CoM
offset only through \eqref{eq:tcom}.

\subsection{Payload-induced parameters}\label{sec:model:payload}

We model the payload as a point mass $\mpay$ rigidly attached at a known
body-frame offset $\rp = [r_x, r_y, r_z]^{\!\top}$ (with $r_z < 0$), and
write the composite mass and the reduced mass as
\begin{equation}
\mt = \mb + \mpay, \qquad \mu = \frac{\mb\, \mpay}{\mt}.
\label{eq:mu}
\end{equation}

\paragraph{CoM offset} The composite CoM is the mass-weighted mean
position, $\rc = (\mpay/\mt)\, \rp$, whence
\begin{equation}
\cvec = [c_x,\, c_y] = \frac{\mpay}{\mt}\, [r_x,\, r_y].
\label{eq:cxy}
\end{equation}

\paragraph{Inertia increment} Because the CoM \emph{moves} when the
payload is attached, the parallel-axis theorem must be applied about the
new CoM, and \emph{both} bodies contribute: the airframe now sits at
distance $(\mpay/\mt)\|\rp\|$ from it and the payload at
$(\mb/\mt)\|\rp\|$. For a lever arm $d$ the two terms collapse onto the
reduced mass,
\begin{equation}
\mb \Bigl(\frac{\mpay}{\mt}\Bigr)^{\!2} d^2 + \mpay \Bigl(\frac{\mb}{\mt}\Bigr)^{\!2} d^2
= \frac{\mb \mpay}{\mt}\, d^2 = \mu\, d^2 ,
\label{eq:reduced}
\end{equation}
which is why $\mu$ and not $\mpay$ is the correct coefficient. Writing
$\mpay d^2$ instead --- the natural slip, treating the payload as orbiting
a fixed body origin --- overestimates the increment by the factor
$\mt/\mb$, here \SI{14.5}{\percent}. Applying \eqref{eq:reduced} per axis
and taking the roll/pitch mean gives the single scalar the model carries,
\begin{align}
\dJ_{xx} &= \mu\bigl(r_y^2 + r_z^2\bigr), \quad
\dJ_{yy} = \mu\bigl(r_x^2 + r_z^2\bigr), \nonumber\\
\dJ &= \tfrac{1}{2}\bigl(\dJ_{xx} + \dJ_{yy}\bigr) = \mu\Bigl( r_z^2 + \tfrac{1}{2}\bigl(r_x^2 + r_y^2\bigr) \Bigr).
\label{eq:dj}
\end{align}

\paragraph{One unknown, three numbers} Equations
\eqref{eq:cxy}--\eqref{eq:dj} produce all three payload-dependent
quantities from a \emph{single} unknown scalar $\mpay$ together with known
geometry $\rp$. The rigid-body constraint thus replaces two degrees of
freedom that a naive augmentation would have to identify from data; the
estimator is not asked to rediscover the parallel-axis theorem from noise.
Table~\ref{tab:params} collects the resulting identities.

\paragraph{Magnitudes} For a \SI{0.3}{\kg} payload at
$\rp = [0,\, 0.109,\, -0.47]\,$\si{\meter} on our platform
($\mb = \SI{2.0643}{\kg}$, $J_{xx} = \SI{0.0142}{\kg\meter\squared}$),
\eqref{eq:mu}--\eqref{eq:dj} give $\mu = \SI{0.262}{\kg}$ and
$\dJ = \SI{0.0594}{\kg\meter\squared}$, an inertia ratio
$(J_{xx}+\dJ)/J_{xx} = 5.18$, together with $c_y = \SI{13.8}{\milli\meter}$
and hence a constant hover trim torque $|c_y|\mt g = \SI{0.32}{\newton\meter}$,
about \SI{64}{\percent} of the roll authority. A \SI{0.3}{\kg} parcel is
therefore not a \SI{15}{\percent} perturbation of the mass: it is a
five-fold change of inertia plus a standing demand on two thirds of the
torque margin. \emph{Adapting the mass alone cannot account for either.}

\begin{table*}[t]
\caption{Model parameters: which are estimated online, which are algebraically
derived, and which are fixed. \emph{Only the mass is a decision variable of the
estimator}; the inertia increment and CoM offset enter the MHE as \emph{known}
runtime parameters, on the same footing as the measured thrust and torque.}
\label{tab:params}
\centering
\footnotesize
\begin{tabular}{@{}l p{0.17\textwidth} c p{0.25\textwidth} p{0.28\textwidth}@{}}
\toprule
Symbol & Meaning & Dim. & Identity & Enters the model through \\
\midrule
$m$ & composite mass (airframe + payload) & 1 & \textbf{MHE decision variable} ($\dot m = 0$, state 14) & translational dynamics \eqref{eq:trans}; cost hover baseline \eqref{eq:hover} \\
$\dJ$ & roll/pitch inertia increment & 1 & algebraically derived, \eqref{eq:dj} & composite inertia \eqref{eq:Jt} in \eqref{eq:rot}; gain scaling \eqref{eq:gain} \\
$\cvec = [c_x, c_y]$ & composite-CoM horizontal offset & 2 & derived \eqref{eq:cxy}, or algebraically inverted online \eqref{eq:conline} & thrust eccentricity torque \eqref{eq:tcom}; cost torque baseline \eqref{eq:hover} \\
\midrule
$\mpay$ & payload-mass scalar (\emph{the single unknown behind $\dJ$ and $\cvec$}) & 1 & weak operational prior (nominal parcel mass) & \eqref{eq:cxy}, \eqref{eq:dj} \\
$\rp = [r_x, r_y, r_z]$ & payload offset in body frame & 3 & known (gripper contact geometry) & \eqref{eq:cxy}, \eqref{eq:dj} \\
\midrule
$\mb$ & bare-airframe mass & 1 & fixed calibration & \SI{2.0643}{\kg} \\
$\Jb$ & bare-airframe inertia & 3 & fixed calibration & $\mathrm{diag}(0.0142, 0.0142, 0.0210)$ \si{\kg\meter\squared} \\
$k_d$, $g$ & drag coefficient, gravity & 2 & fixed constants & $0.05$, \SI{9.81}{\meter\per\second\squared} \\
\bottomrule
\end{tabular}
\end{table*}

\subsection{Model-independent reconstruction of the physical input}\label{sec:model:phys}

Both $T$ and $\boldsymbol{\tau}$ in \eqref{eq:trans}--\eqref{eq:rot} are
reconstructed from measured rotor speeds $\omega_i$ and the known rotor
positions $(x_i, y_i)$, with $F_i = k_f \omega_i^2$,
\begin{equation}
\Tphys = \sum_i F_i, \quad
\tau_{\mathrm{roll}} = \sum_i y_i F_i, \quad
\tau_{\mathrm{pitch}} = -\sum_i x_i F_i .
\label{eq:phys}
\end{equation}
Feeding the controller's \emph{commanded} thrust back into the estimator
instead would create a self-consistent blind spot: at hover
$m \approx T/g$ holds for \emph{any} $m$ if $T$ is itself computed from the
assumed $m$, and the mass becomes unobservable. Equation \eqref{eq:phys}
breaks the loop because it reads a physical signal that no model state
enters.

\subsection{Estimation: the mass as the single online degree of freedom}\label{sec:model:mhe}

The moving-horizon estimator augments the mass into the state,
\begin{equation}
x_{\mathrm{aug}} = [x;\, m] \in \mathbb{R}^{14}, \qquad \dot m = 0,
\label{eq:xaug}
\end{equation}
holding it constant over the window so that the horizon of measurements
constrains one value rather than a free trajectory. Process noise
$\mathbf{w} \in \mathbb{R}^{13}$ acts on the physical states only; adding
noise on the mass channel would let it drift within the window and dissolve
the identification constraint. Over a window of $N$ stages the estimator
solves
\begin{align}
\min_{x_{0}, \mathbf{w}} \quad
& \sum_{k=0}^{N-1} \Bigl( \| \mathbf{y}_k - h(x_k) \|^2_{R} + \| \mathbf{w}_k \|^2_{Q} \Bigr)
+ \| x_0 - \bar{x}_0 \|^2_{Q_0} \nonumber \\
\text{s.t.} \quad & x_{k+1} = f\bigl(x_k, \mathbf{w}_k;\, p_{\mathrm{MHE}}\bigr), \quad
m \in [\,\underline{m},\, \overline{m}\,],
\label{eq:mhe}
\end{align}
with $f$ the discretisation of \eqref{eq:kin}--\eqref{eq:rot}, arrival cost
$\| \cdot \|^2_{Q_0}$ anchoring the window start, and measurement
$\mathbf{y}$ the 13 physical states. Crucially, the payload geometry is
\emph{not} in the decision vector but in the runtime parameter vector,
alongside the measured input:
\begin{equation}
p_{\mathrm{MHE}} = \bigl[\, T,\; \boldsymbol{\tau},\; \dJ,\; c_x,\; c_y \,\bigr] \in \mathbb{R}^{7}.
\label{eq:pmhe}
\end{equation}
Inside the estimator, $\dJ$ and $\cvec$ are \emph{known inputs}, on the
same footing as $T$ and $\boldsymbol{\tau}$. We use $N = 20$ at
$dt = \SI{0.1}{\second}$ (a \SI{2.0}{\second} horizon) and bound the mass
to $[\SI{1.961}{\kg}, \SI{5.0}{\kg}]$; the lower bound is the physical
floor $0.95\,\mb$, since the vehicle cannot weigh less than its own
airframe. The bound is a hard constraint on the decision variable --- an
advantage of the moving-horizon formulation over a recursive filter, which
can only be nudged away from non-physical estimates.

\paragraph{Ground-truth-free CoM inversion} When rotor speeds are
available the CoM offset need not come from the prior at all. In steady
hover the eccentric moment \eqref{eq:tcom} is balanced by the inner loop,
so $\dot{\boldsymbol{\omega}} \approx 0$ and \eqref{eq:rot} inverts
directly:
\begin{equation}
c_x = -\,\tau_{\mathrm{pitch}}^{\mathrm{phys}} / \Tphys, \qquad
c_y = +\,\tau_{\mathrm{roll}}^{\mathrm{phys}} / \Tphys .
\label{eq:conline}
\end{equation}
This is an algebraic inversion, not an optimisation, and it adds no degree
of freedom. Three properties make it admissible. \emph{(i)} It reads no
model state, so it cannot close an estimate-to-model feedback loop.
\emph{(ii)} The rotor coefficient $k_f$ cancels in the ratio ---
$c_x = \sum_i x_i \omega_i^2 / \sum_i \omega_i^2$ --- so the result is
immune to thrust-calibration error, unlike any scheme that consumes an
absolute thrust. \emph{(iii)} Channels are assigned by identifiability:
$\cvec$ has an independent observable and is estimated online, whereas
$r_z$ is annihilated by the cross product in \eqref{eq:tcom} and cannot be
identified from the torque channel at all --- it must remain a prior.
Samples are gated to steady flight
($\|\boldsymbol{\omega}\| < \SI{0.15}{\radian\per\second}$,
$\|\mathbf{v}_{xy}\| < \SI{0.20}{\meter\per\second}$), which rejects the
\emph{systematic} error of manoeuvring torques, and then low-pass filtered
with an exponential moving average of time constant \SI{2}{\second}, which
suppresses \emph{random} error: $\cvec$ is quasi-static once the load is
grasped, so its signal lies at DC while the differential-mode rotor noise
does not.

\subsection{Control: how the parameters re-centre the cost}\label{sec:model:ctrl}

The NMPC carries the parameters as runtime values,
\begin{equation}
p_{\mathrm{NMPC}} = \bigl[\, x_{\mathrm{ref}}(13),\; m,\; \dJ,\; \cvec(2),\; \mathbf{d}(3) \,\bigr] \in \mathbb{R}^{20},
\label{eq:pnmpc}
\end{equation}
where $\mathbf{d}$ is a lumped translational disturbance used only by the
$\mathcal{L}_1$ baseline of \S\ref{sec:results}-E and identically zero
under the proposed method. It minimises, over a horizon of $N$ stages,
\begin{equation}
\sum_{k} \Bigl( \| e_{\mathrm{track}}(x_k, x_{\mathrm{ref},k}) \|^2_{Q}
+ \| u_k - u_{\mathrm{hover}} \|^2_{R} \Bigr)
\label{eq:nmpc}
\end{equation}
subject to \eqref{eq:kin}--\eqref{eq:rot} and to input bounds on thrust and
torque. The second term penalises the deviation of the input from a
\emph{baseline}, and that baseline is recomputed from the live parameters
rather than fixed at build time:
\begin{equation}
T_h = m\,g, \qquad
u_{\mathrm{hover}} = \bigl[\, T_h,\; c_y T_h,\; -c_x T_h,\; 0 \,\bigr]^{\!\top}.
\label{eq:hover}
\end{equation}
Both entries matter, and for the same reason. Penalising $\|u\|^2_R$
outright would treat the thrust needed to hover as an expense to be
minimised; anchoring instead at a \emph{stale} $T_h$ merely moves the
error --- with a baseline computed from the unloaded mass, the cost pulls
the thrust down persistently and the tracking error settles at a nonzero
offset that no amount of retuning removes (increasing $R$ makes it worse,
since $R$ is the strength with which the wrong baseline pulls). The torque
entries are exactly $-\tcom$ of \eqref{eq:tcom}: hovering with an eccentric
load requires a constant trim torque, so a cost that does not admit it
pulls the torque toward zero and produces the same class of steady-state
bias. Equation \eqref{eq:tcom} teaches the \emph{model} that the moment
exists; equation \eqref{eq:hover} teaches the \emph{objective} that it is
necessary. Both are required.

\paragraph{Inner-loop bandwidth} Model consistency in the outer loop
cannot restore inner-loop bandwidth. A payload that multiplies roll/pitch
inertia lowers the effective bandwidth of the PX4 rate controller, and for
heavy eccentric loads the resulting attitude oscillation grows into torque
saturation. We therefore scale the rate gains at attach by the inertia
ratio,
\begin{equation}
K \leftarrow K_{\mathrm{nom}} \cdot \min\!\Bigl( \frac{J_{xx} + \dJ}{J_{xx}},\; 5.0 \Bigr).
\label{eq:gain}
\end{equation}
The cap is not incidental: at \SI{0.3}{\kg} the ratio is already $5.18$,
so that payload sits at the ceiling of what inner-loop retuning alone can
absorb on this airframe --- a boundary we quantify in
\S\ref{sec:results}-C.

\subsection{Why the inertia and CoM are not augmented into the state}\label{sec:model:whynot}

Augmenting $\dJ$ and $\cvec$ into the MHE decision vector is the obvious
alternative. We argue against it on four grounds.

\emph{(i) Asymmetric observability.} The mass is excited directly by the
thrust--acceleration relation \eqref{eq:trans} and is observable in
near-hover flight. The inertia increment acts only through
$\dot{\boldsymbol{\omega}}$ in \eqref{eq:rot}, which a delivery task ---
spent largely in near-hover, low-rate flight --- excites weakly. Making it
free lets the optimiser explain measurement noise with a nearly
unobservable quantity.

\emph{(ii) They compete for one residual.} Deriving the geometry from the
\emph{instantaneous} mass estimate closes a feedback loop
$m_{\mathrm{est}} \rightarrow \dJ, \cvec \rightarrow \text{attitude model}
\rightarrow m_{\mathrm{est}}$. We observed this bootstrap fail structurally
on light payloads: before lift-off the mass estimate reads low
\emph{legitimately}, that value is consumed as ``no geometry'', the
geometry channel closes, and after lift-off the model is mismatched and
the estimate sticks at its lower bound (\S\ref{sec:method}-E). The fix was
to sever the loop, not to add more free parameters.

\emph{(iii) Economy of parameterisation.} Equations
\eqref{eq:cxy}--\eqref{eq:dj} obtain three numbers from one scalar. Two
degrees of freedom are bought back by a constraint that holds exactly for a
rigid attachment.

\emph{(iv) The remaining prior may be weak.} With the geometry prior set
\SI{33}{\percent} wrong, the mass estimate still converges to within
\SI{0.3}{\percent} (\S\ref{sec:results}-C), so a nominal parcel mass
suffices on hardware and no attachment ground truth is required. Where
rotor speeds are available, \eqref{eq:conline} removes the prior from the
$\cvec$ path entirely.

\subsection{Declared modeling approximations}\label{sec:model:approx}

Three points deserve explicit statement rather than silence.

\emph{(a) Anisotropy of $\dJ$.} The single scalar in \eqref{eq:dj}
replaces the pair $(\dJ_{xx}, \dJ_{yy})$ by their mean, an error of
$\tfrac{1}{2}\mu|r_x^2 - r_y^2|$, below \SI{1}{\percent} of $J_{xx}+\dJ$ at
our offsets, and exactly zero when $r_x = r_y$.

\emph{(b) The yaw increment is neglected.} $\dJ_{zz} = \mu(r_x^2 + r_y^2)$
is not added to $J_{zz}$ in \eqref{eq:Jt}. At \SI{0.3}{\kg} and
$r_{xy} = \SI{0.109}{\meter}$ it amounts to \SI{0.0031}{\kg\meter\squared},
some \SI{15}{\percent} of $J_{zz}$ --- a known, non-negligible
approximation. It does not affect the reported results, since the yaw
channel neither carries payload compensation nor participates in
\eqref{eq:conline}; and removing it requires only substituting
$J_{zz} + \mu(r_x^2 + r_y^2)$ in \eqref{eq:Jt}, introducing no new
estimated quantity.

\emph{(c) $c_z$ is absent by construction.} As shown below
\eqref{eq:tcom}, this is an exact algebraic consequence, not a truncation,
and is listed here only to forestall the opposite reading.
   (place after \section{Related Work};
%%         it supersedes the "Dynamics and parameterization" / "MHE" /
%%         "Inner-loop gain scaling" paragraphs of the current
%%         \section{System Overview and Problem Setup}, and subsumes
%%         sec3_parameters.tex in full.)
%%
%% REQUIRES in the preamble (all already present in main.tex):
%%   amsmath, amssymb, booktabs, siunitx, multirow, array
%% =====================================================================

% ---- Notation macros (safe to move to the preamble) -----------------
\providecommand{\dJ}{\Delta J}
\providecommand{\Jt}{J_t}                  % composite inertia
\providecommand{\Jb}{J_b}                  % bare-airframe inertia
\providecommand{\cvec}{\mathbf{c}}
\providecommand{\rp}{\mathbf{r}_p}
\providecommand{\rc}{\mathbf{r}_c}
\providecommand{\mpay}{\hat{m}_p}          % payload-mass scalar (prior)
\providecommand{\mb}{m_b}                  % bare-airframe mass
\providecommand{\mt}{m_t}                  % composite mass
\providecommand{\Tphys}{T_{\mathrm{phys}}}
\providecommand{\tauphys}{\boldsymbol{\tau}_{\mathrm{phys}}}
\providecommand{\tcom}{\boldsymbol{\tau}_{\mathrm{com}}}
\providecommand{\ev}{\mathbf{e}_3}

\section{Modeling and Parameter Structure}\label{sec:model}

This section states the model the estimator and the controller share, and
then makes explicit \emph{which} of its parameters are estimated online.
The distinction is the substance of the method: a payload changes three
quantities in the dynamics --- mass, roll/pitch inertia, and the
horizontal offset of the composite centre of mass (CoM) --- yet only one
of them is a decision variable of the estimator. The remaining two are
generated from it and from known contact geometry by rigid-body algebra.

\subsection{Frames, notation, and conventions}\label{sec:model:frames}

We use an ENU world frame $\mathcal{W}$ and an FLU body frame
$\mathcal{B}$ whose origin is the geometric centre of the rotor disc.
The state and input are
\begin{equation}
x = \bigl[\mathbf{p};\, \mathbf{v};\, \mathbf{q};\, \boldsymbol{\omega}\bigr] \in \mathbb{R}^{13},
\qquad
u = \bigl[T;\, \boldsymbol{\tau}\bigr] \in \mathbb{R}^{4},
\label{eq:state}
\end{equation}
with position $\mathbf{p}$ and velocity $\mathbf{v}$ in $\mathcal{W}$,
body rate $\boldsymbol{\omega}$ in $\mathcal{B}$, collective thrust $T$
along the body $z$ axis, and body torque $\boldsymbol{\tau}$. We write
$\ev = [0,0,1]^{\!\top}$.

Attitude is a unit quaternion under the \emph{Hamilton} convention,
scalar first, $\mathbf{q} = [q_w, q_x, q_y, q_z]^{\!\top}$, representing
the rotation from $\mathcal{B}$ to $\mathcal{W}$; $R(\mathbf{q}) \in SO(3)$
denotes the corresponding rotation matrix. Stating the convention matters
because the body-rate kinematics \eqref{eq:quat} is a \emph{right}
multiplication precisely when $\boldsymbol{\omega}$ is resolved in
$\mathcal{B}$ --- the JPL convention, or a world-frame rate, would place
the product on the other side and integrate the attitude the wrong way.

\subsection{Rigid-body kinematics}\label{sec:model:kin}

Translational kinematics is trivial,
\begin{equation}
\dot{\mathbf{p}} = \mathbf{v}. \label{eq:kin}
\end{equation}
Attitude propagates by the quaternion product with the pure quaternion
$[0, \boldsymbol{\omega}]$,
\begin{equation}
\dot{\mathbf{q}} = \tfrac{1}{2}\, \mathbf{q} \otimes [0,\, \boldsymbol{\omega}]
= \tfrac{1}{2}
\begin{bmatrix}
-q_x & -q_y & -q_z \\
 q_w & -q_z &  q_y \\
 q_z &  q_w & -q_x \\
-q_y &  q_x &  q_w
\end{bmatrix} \boldsymbol{\omega},
\label{eq:quat}
\end{equation}
the two forms being algebraically identical; the implementation uses the
$4 \times 3$ matrix on the right. Neither \eqref{eq:kin} nor
\eqref{eq:quat} contains a payload-dependent parameter: \emph{the payload
enters the model only through the dynamics}, which is what confines the
parameter discussion below to three equations.

\subsection{Rigid-body dynamics}\label{sec:model:dyn}

\paragraph{Translation} Newton's law in $\mathcal{W}$, with a linear
aerodynamic drag term, reads
\begin{equation}
\dot{\mathbf{v}} = \frac{1}{m}\Bigl( R(\mathbf{q})\, \ev\, T - k_d\, \mathbf{v} \Bigr) - g\, \ev,
\label{eq:trans}
\end{equation}
where $m$ is the \emph{composite} mass (airframe plus payload) and $k_d$ a
fixed drag coefficient. The mass appears in exactly one role, as the
thrust-to-acceleration gain $1/m$. This is the mechanism of its
observability: given a thrust value reconstructed independently of the
model (\S\ref{sec:model:phys}), the measured $\dot{\mathbf{v}}$ determines
$m$ uniquely, and it does so in near-hover flight, without a dedicated
excitation manoeuvre.

\paragraph{Rotation} Let $\Jt$ denote the composite inertia, i.e.\ the
bare-airframe value $\Jb = \mathrm{diag}(J_{xx}, J_{yy}, J_{zz})$ raised by
the payload increment $\dJ$ on the roll and pitch axes,
\begin{equation}
\Jt = \mathrm{diag}\bigl(J_{xx} + \dJ,\; J_{yy} + \dJ,\; J_{zz}\bigr),
\label{eq:Jt}
\end{equation}
in the same spirit as the composite mass $\mt$ of
\S\ref{sec:model:payload}. Euler's equation about the composite CoM is
then
\begin{equation}
\dot{\boldsymbol{\omega}} = \Jt^{-1}\Bigl( \boldsymbol{\tau} + \tcom(\cvec) - \boldsymbol{\omega} \times \Jt \boldsymbol{\omega} \Bigr).
\label{eq:rot}
\end{equation}

\paragraph{Thrust eccentricity} The term $\tcom$ is the moment the
collective thrust exerts about the composite CoM. Thrust acts at the
rotor-disc centre --- the body origin --- along body $z$, while the CoM
sits at $\rc = [c_x, c_y, c_z]^{\!\top}$, so
\begin{equation}
\tcom(\cvec) = (-\rc) \times \ev T = \bigl[-c_y T,\; +c_x T,\; 0\bigr]^{\!\top}.
\label{eq:tcom}
\end{equation}
Two consequences are worth stating explicitly. First, $c_z$ \emph{cannot}
appear: a vertical CoM offset is collinear with the body-$z$ thrust, so
their cross product vanishes identically, and gravity acts at the CoM
itself and exerts no moment about it. The CoM offset is a two-parameter
quantity \emph{by construction, not by truncation} --- a point worth making
because ``the offset is three-dimensional yet only two components are
estimated'' is a natural objection. Second, \eqref{eq:tcom} is a
\emph{constant} moment in level hover, and a large one: at the payload of
\S\ref{sec:model:payload} it consumes roughly two thirds of the available
roll authority. Hovering on the spot with an off-centre load therefore
\emph{requires} a persistent trim torque, a fact the cost function must
also admit (\S\ref{sec:model:ctrl}).

Equations \eqref{eq:trans}--\eqref{eq:tcom} exhibit the property this
paper builds on: \textbf{each payload-dependent parameter has exactly one
point of entry.} The mass enters only through $1/m$ in \eqref{eq:trans},
the inertia increment only through $\Jt$ in \eqref{eq:rot}, and the CoM
offset only through \eqref{eq:tcom}.

\subsection{Payload-induced parameters}\label{sec:model:payload}

We model the payload as a point mass $\mpay$ rigidly attached at a known
body-frame offset $\rp = [r_x, r_y, r_z]^{\!\top}$ (with $r_z < 0$), and
write the composite mass and the reduced mass as
\begin{equation}
\mt = \mb + \mpay, \qquad \mu = \frac{\mb\, \mpay}{\mt}.
\label{eq:mu}
\end{equation}

\paragraph{CoM offset} The composite CoM is the mass-weighted mean
position, $\rc = (\mpay/\mt)\, \rp$, whence
\begin{equation}
\cvec = [c_x,\, c_y] = \frac{\mpay}{\mt}\, [r_x,\, r_y].
\label{eq:cxy}
\end{equation}

\paragraph{Inertia increment} Because the CoM \emph{moves} when the
payload is attached, the parallel-axis theorem must be applied about the
new CoM, and \emph{both} bodies contribute: the airframe now sits at
distance $(\mpay/\mt)\|\rp\|$ from it and the payload at
$(\mb/\mt)\|\rp\|$. For a lever arm $d$ the two terms collapse onto the
reduced mass,
\begin{equation}
\mb \Bigl(\frac{\mpay}{\mt}\Bigr)^{\!2} d^2 + \mpay \Bigl(\frac{\mb}{\mt}\Bigr)^{\!2} d^2
= \frac{\mb \mpay}{\mt}\, d^2 = \mu\, d^2 ,
\label{eq:reduced}
\end{equation}
which is why $\mu$ and not $\mpay$ is the correct coefficient. Writing
$\mpay d^2$ instead --- the natural slip, treating the payload as orbiting
a fixed body origin --- overestimates the increment by the factor
$\mt/\mb$, here \SI{14.5}{\percent}. Applying \eqref{eq:reduced} per axis
and taking the roll/pitch mean gives the single scalar the model carries,
\begin{align}
\dJ_{xx} &= \mu\bigl(r_y^2 + r_z^2\bigr), \quad
\dJ_{yy} = \mu\bigl(r_x^2 + r_z^2\bigr), \nonumber\\
\dJ &= \tfrac{1}{2}\bigl(\dJ_{xx} + \dJ_{yy}\bigr) = \mu\Bigl( r_z^2 + \tfrac{1}{2}\bigl(r_x^2 + r_y^2\bigr) \Bigr).
\label{eq:dj}
\end{align}

\paragraph{One unknown, three numbers} Equations
\eqref{eq:cxy}--\eqref{eq:dj} produce all three payload-dependent
quantities from a \emph{single} unknown scalar $\mpay$ together with known
geometry $\rp$. The rigid-body constraint thus replaces two degrees of
freedom that a naive augmentation would have to identify from data; the
estimator is not asked to rediscover the parallel-axis theorem from noise.
Table~\ref{tab:params} collects the resulting identities.

\paragraph{Magnitudes} For a \SI{0.3}{\kg} payload at
$\rp = [0,\, 0.109,\, -0.47]\,$\si{\meter} on our platform
($\mb = \SI{2.0643}{\kg}$, $J_{xx} = \SI{0.0142}{\kg\meter\squared}$),
\eqref{eq:mu}--\eqref{eq:dj} give $\mu = \SI{0.262}{\kg}$ and
$\dJ = \SI{0.0594}{\kg\meter\squared}$, an inertia ratio
$(J_{xx}+\dJ)/J_{xx} = 5.18$, together with $c_y = \SI{13.8}{\milli\meter}$
and hence a constant hover trim torque $|c_y|\mt g = \SI{0.32}{\newton\meter}$,
about \SI{64}{\percent} of the roll authority. A \SI{0.3}{\kg} parcel is
therefore not a \SI{15}{\percent} perturbation of the mass: it is a
five-fold change of inertia plus a standing demand on two thirds of the
torque margin. \emph{Adapting the mass alone cannot account for either.}

\begin{table*}[t]
\caption{Model parameters: which are estimated online, which are algebraically
derived, and which are fixed. \emph{Only the mass is a decision variable of the
estimator}; the inertia increment and CoM offset enter the MHE as \emph{known}
runtime parameters, on the same footing as the measured thrust and torque.}
\label{tab:params}
\centering
\footnotesize
\begin{tabular}{@{}l p{0.17\textwidth} c p{0.25\textwidth} p{0.28\textwidth}@{}}
\toprule
Symbol & Meaning & Dim. & Identity & Enters the model through \\
\midrule
$m$ & composite mass (airframe + payload) & 1 & \textbf{MHE decision variable} ($\dot m = 0$, state 14) & translational dynamics \eqref{eq:trans}; cost hover baseline \eqref{eq:hover} \\
$\dJ$ & roll/pitch inertia increment & 1 & algebraically derived, \eqref{eq:dj} & composite inertia \eqref{eq:Jt} in \eqref{eq:rot}; gain scaling \eqref{eq:gain} \\
$\cvec = [c_x, c_y]$ & composite-CoM horizontal offset & 2 & derived \eqref{eq:cxy}, or algebraically inverted online \eqref{eq:conline} & thrust eccentricity torque \eqref{eq:tcom}; cost torque baseline \eqref{eq:hover} \\
\midrule
$\mpay$ & payload-mass scalar (\emph{the single unknown behind $\dJ$ and $\cvec$}) & 1 & weak operational prior (nominal parcel mass) & \eqref{eq:cxy}, \eqref{eq:dj} \\
$\rp = [r_x, r_y, r_z]$ & payload offset in body frame & 3 & known (gripper contact geometry) & \eqref{eq:cxy}, \eqref{eq:dj} \\
\midrule
$\mb$ & bare-airframe mass & 1 & fixed calibration & \SI{2.0643}{\kg} \\
$\Jb$ & bare-airframe inertia & 3 & fixed calibration & $\mathrm{diag}(0.0142, 0.0142, 0.0210)$ \si{\kg\meter\squared} \\
$k_d$, $g$ & drag coefficient, gravity & 2 & fixed constants & $0.05$, \SI{9.81}{\meter\per\second\squared} \\
\bottomrule
\end{tabular}
\end{table*}

\subsection{Model-independent reconstruction of the physical input}\label{sec:model:phys}

Both $T$ and $\boldsymbol{\tau}$ in \eqref{eq:trans}--\eqref{eq:rot} are
reconstructed from measured rotor speeds $\omega_i$ and the known rotor
positions $(x_i, y_i)$, with $F_i = k_f \omega_i^2$,
\begin{equation}
\Tphys = \sum_i F_i, \quad
\tau_{\mathrm{roll}} = \sum_i y_i F_i, \quad
\tau_{\mathrm{pitch}} = -\sum_i x_i F_i .
\label{eq:phys}
\end{equation}
Feeding the controller's \emph{commanded} thrust back into the estimator
instead would create a self-consistent blind spot: at hover
$m \approx T/g$ holds for \emph{any} $m$ if $T$ is itself computed from the
assumed $m$, and the mass becomes unobservable. Equation \eqref{eq:phys}
breaks the loop because it reads a physical signal that no model state
enters.

\subsection{Estimation: the mass as the single online degree of freedom}\label{sec:model:mhe}

The moving-horizon estimator augments the mass into the state,
\begin{equation}
x_{\mathrm{aug}} = [x;\, m] \in \mathbb{R}^{14}, \qquad \dot m = 0,
\label{eq:xaug}
\end{equation}
holding it constant over the window so that the horizon of measurements
constrains one value rather than a free trajectory. Process noise
$\mathbf{w} \in \mathbb{R}^{13}$ acts on the physical states only; adding
noise on the mass channel would let it drift within the window and dissolve
the identification constraint. Over a window of $N$ stages the estimator
solves
\begin{align}
\min_{x_{0}, \mathbf{w}} \quad
& \sum_{k=0}^{N-1} \Bigl( \| \mathbf{y}_k - h(x_k) \|^2_{R} + \| \mathbf{w}_k \|^2_{Q} \Bigr)
+ \| x_0 - \bar{x}_0 \|^2_{Q_0} \nonumber \\
\text{s.t.} \quad & x_{k+1} = f\bigl(x_k, \mathbf{w}_k;\, p_{\mathrm{MHE}}\bigr), \quad
m \in [\,\underline{m},\, \overline{m}\,],
\label{eq:mhe}
\end{align}
with $f$ the discretisation of \eqref{eq:kin}--\eqref{eq:rot}, arrival cost
$\| \cdot \|^2_{Q_0}$ anchoring the window start, and measurement
$\mathbf{y}$ the 13 physical states. Crucially, the payload geometry is
\emph{not} in the decision vector but in the runtime parameter vector,
alongside the measured input:
\begin{equation}
p_{\mathrm{MHE}} = \bigl[\, T,\; \boldsymbol{\tau},\; \dJ,\; c_x,\; c_y \,\bigr] \in \mathbb{R}^{7}.
\label{eq:pmhe}
\end{equation}
Inside the estimator, $\dJ$ and $\cvec$ are \emph{known inputs}, on the
same footing as $T$ and $\boldsymbol{\tau}$. We use $N = 20$ at
$dt = \SI{0.1}{\second}$ (a \SI{2.0}{\second} horizon) and bound the mass
to $[\SI{1.961}{\kg}, \SI{5.0}{\kg}]$; the lower bound is the physical
floor $0.95\,\mb$, since the vehicle cannot weigh less than its own
airframe. The bound is a hard constraint on the decision variable --- an
advantage of the moving-horizon formulation over a recursive filter, which
can only be nudged away from non-physical estimates.

\paragraph{Ground-truth-free CoM inversion} When rotor speeds are
available the CoM offset need not come from the prior at all. In steady
hover the eccentric moment \eqref{eq:tcom} is balanced by the inner loop,
so $\dot{\boldsymbol{\omega}} \approx 0$ and \eqref{eq:rot} inverts
directly:
\begin{equation}
c_x = -\,\tau_{\mathrm{pitch}}^{\mathrm{phys}} / \Tphys, \qquad
c_y = +\,\tau_{\mathrm{roll}}^{\mathrm{phys}} / \Tphys .
\label{eq:conline}
\end{equation}
This is an algebraic inversion, not an optimisation, and it adds no degree
of freedom. Three properties make it admissible. \emph{(i)} It reads no
model state, so it cannot close an estimate-to-model feedback loop.
\emph{(ii)} The rotor coefficient $k_f$ cancels in the ratio ---
$c_x = \sum_i x_i \omega_i^2 / \sum_i \omega_i^2$ --- so the result is
immune to thrust-calibration error, unlike any scheme that consumes an
absolute thrust. \emph{(iii)} Channels are assigned by identifiability:
$\cvec$ has an independent observable and is estimated online, whereas
$r_z$ is annihilated by the cross product in \eqref{eq:tcom} and cannot be
identified from the torque channel at all --- it must remain a prior.
Samples are gated to steady flight
($\|\boldsymbol{\omega}\| < \SI{0.15}{\radian\per\second}$,
$\|\mathbf{v}_{xy}\| < \SI{0.20}{\meter\per\second}$), which rejects the
\emph{systematic} error of manoeuvring torques, and then low-pass filtered
with an exponential moving average of time constant \SI{2}{\second}, which
suppresses \emph{random} error: $\cvec$ is quasi-static once the load is
grasped, so its signal lies at DC while the differential-mode rotor noise
does not.

\subsection{Control: how the parameters re-centre the cost}\label{sec:model:ctrl}

The NMPC carries the parameters as runtime values,
\begin{equation}
p_{\mathrm{NMPC}} = \bigl[\, x_{\mathrm{ref}}(13),\; m,\; \dJ,\; \cvec(2),\; \mathbf{d}(3) \,\bigr] \in \mathbb{R}^{20},
\label{eq:pnmpc}
\end{equation}
where $\mathbf{d}$ is a lumped translational disturbance used only by the
$\mathcal{L}_1$ baseline of \S\ref{sec:results}-E and identically zero
under the proposed method. It minimises, over a horizon of $N$ stages,
\begin{equation}
\sum_{k} \Bigl( \| e_{\mathrm{track}}(x_k, x_{\mathrm{ref},k}) \|^2_{Q}
+ \| u_k - u_{\mathrm{hover}} \|^2_{R} \Bigr)
\label{eq:nmpc}
\end{equation}
subject to \eqref{eq:kin}--\eqref{eq:rot} and to input bounds on thrust and
torque. The second term penalises the deviation of the input from a
\emph{baseline}, and that baseline is recomputed from the live parameters
rather than fixed at build time:
\begin{equation}
T_h = m\,g, \qquad
u_{\mathrm{hover}} = \bigl[\, T_h,\; c_y T_h,\; -c_x T_h,\; 0 \,\bigr]^{\!\top}.
\label{eq:hover}
\end{equation}
Both entries matter, and for the same reason. Penalising $\|u\|^2_R$
outright would treat the thrust needed to hover as an expense to be
minimised; anchoring instead at a \emph{stale} $T_h$ merely moves the
error --- with a baseline computed from the unloaded mass, the cost pulls
the thrust down persistently and the tracking error settles at a nonzero
offset that no amount of retuning removes (increasing $R$ makes it worse,
since $R$ is the strength with which the wrong baseline pulls). The torque
entries are exactly $-\tcom$ of \eqref{eq:tcom}: hovering with an eccentric
load requires a constant trim torque, so a cost that does not admit it
pulls the torque toward zero and produces the same class of steady-state
bias. Equation \eqref{eq:tcom} teaches the \emph{model} that the moment
exists; equation \eqref{eq:hover} teaches the \emph{objective} that it is
necessary. Both are required.

\paragraph{Inner-loop bandwidth} Model consistency in the outer loop
cannot restore inner-loop bandwidth. A payload that multiplies roll/pitch
inertia lowers the effective bandwidth of the PX4 rate controller, and for
heavy eccentric loads the resulting attitude oscillation grows into torque
saturation. We therefore scale the rate gains at attach by the inertia
ratio,
\begin{equation}
K \leftarrow K_{\mathrm{nom}} \cdot \min\!\Bigl( \frac{J_{xx} + \dJ}{J_{xx}},\; 5.0 \Bigr).
\label{eq:gain}
\end{equation}
The cap is not incidental: at \SI{0.3}{\kg} the ratio is already $5.18$,
so that payload sits at the ceiling of what inner-loop retuning alone can
absorb on this airframe --- a boundary we quantify in
\S\ref{sec:results}-C.

\subsection{Why the inertia and CoM are not augmented into the state}\label{sec:model:whynot}

Augmenting $\dJ$ and $\cvec$ into the MHE decision vector is the obvious
alternative. We argue against it on four grounds.

\emph{(i) Asymmetric observability.} The mass is excited directly by the
thrust--acceleration relation \eqref{eq:trans} and is observable in
near-hover flight. The inertia increment acts only through
$\dot{\boldsymbol{\omega}}$ in \eqref{eq:rot}, which a delivery task ---
spent largely in near-hover, low-rate flight --- excites weakly. Making it
free lets the optimiser explain measurement noise with a nearly
unobservable quantity.

\emph{(ii) They compete for one residual.} Deriving the geometry from the
\emph{instantaneous} mass estimate closes a feedback loop
$m_{\mathrm{est}} \rightarrow \dJ, \cvec \rightarrow \text{attitude model}
\rightarrow m_{\mathrm{est}}$. We observed this bootstrap fail structurally
on light payloads: before lift-off the mass estimate reads low
\emph{legitimately}, that value is consumed as ``no geometry'', the
geometry channel closes, and after lift-off the model is mismatched and
the estimate sticks at its lower bound (\S\ref{sec:method}-E). The fix was
to sever the loop, not to add more free parameters.

\emph{(iii) Economy of parameterisation.} Equations
\eqref{eq:cxy}--\eqref{eq:dj} obtain three numbers from one scalar. Two
degrees of freedom are bought back by a constraint that holds exactly for a
rigid attachment.

\emph{(iv) The remaining prior may be weak.} With the geometry prior set
\SI{33}{\percent} wrong, the mass estimate still converges to within
\SI{0.3}{\percent} (\S\ref{sec:results}-C), so a nominal parcel mass
suffices on hardware and no attachment ground truth is required. Where
rotor speeds are available, \eqref{eq:conline} removes the prior from the
$\cvec$ path entirely.

\subsection{Declared modeling approximations}\label{sec:model:approx}

Three points deserve explicit statement rather than silence.

\emph{(a) Anisotropy of $\dJ$.} The single scalar in \eqref{eq:dj}
replaces the pair $(\dJ_{xx}, \dJ_{yy})$ by their mean, an error of
$\tfrac{1}{2}\mu|r_x^2 - r_y^2|$, below \SI{1}{\percent} of $J_{xx}+\dJ$ at
our offsets, and exactly zero when $r_x = r_y$.

\emph{(b) The yaw increment is neglected.} $\dJ_{zz} = \mu(r_x^2 + r_y^2)$
is not added to $J_{zz}$ in \eqref{eq:Jt}. At \SI{0.3}{\kg} and
$r_{xy} = \SI{0.109}{\meter}$ it amounts to \SI{0.0031}{\kg\meter\squared},
some \SI{15}{\percent} of $J_{zz}$ --- a known, non-negligible
approximation. It does not affect the reported results, since the yaw
channel neither carries payload compensation nor participates in
\eqref{eq:conline}; and removing it requires only substituting
$J_{zz} + \mu(r_x^2 + r_y^2)$ in \eqref{eq:Jt}, introducing no new
estimated quantity.

\emph{(c) $c_z$ is absent by construction.} As shown below
\eqref{eq:tcom}, this is an exact algebraic consequence, not a truncation,
and is listed here only to forestall the opposite reading.
